\documentclass[Master]{subfiles}
\begin{document}
\section{Discussion and summary}
We developed a computationally efficient framework for the computation of second order expectation values and out of time order propagators of second-order Hamiltonians.  We used the framework to visualise the adiabatic exchange of non-abelian Majorana modes in 2D chiral superconductor vortices in real space. In this instance, the, the framework demonstrated promising efficiency for more extensive calculations in future research.\\

The analytical calculations \cite{volovik1999fermion} \cite{ivanov2001non} \cite{stone2004edge} identify circular topologically trivial regions infinitely far apart and of vanishing radius surrounded by a topological chiral$\pm$ bulk with a discrete jump in chemical potential at the edge and a $2\pi$ phase winding of the order parameter around the topological regions as the ideal geometry for non-abelian Majorana modes. \\

However, in a real-world implementation, vortices can neither be infinitely far apart nor have a vanishing radius, and must thus have a small energy splitting. We calculated this energy splitting for vortices with finite distance and radius, and found a further case of exponential energy splitting in dependence from the vortex to the edge of the bulk. The vortex radius was found to have a much smaller effect on the energy splitting than the vortex to vortex/bulk edge distance. Reducing the Vortex radius from $20$ to a single site on a $80\times 80$ grid reduced the energy splitting by one order of magnitude, whereas the energy splitting for vortices touching each other or the bulks edge was seven magnitudes larger that that for maximal distances on a $80\times 80$ grid.
We observe that even vortices as large as $9$ sited across with maximal vortex to vortex/bulk edge distances on a $36\times36$ lattice with discrete edges are sufficient to yield an energy-splitting nine orders of magnitude smaller than the energy-splitting consists of a few of $\mu,t,\Delta$. For a typical hopping term $t$ in the range of few $meV$, the energy-splitting are a few  $peV$. Furthermore, we observed Majorana modes at the interface between the chiral+ and chiral- regime and saw that vortices with a continuous transition of the chemical potential, have larger energy splitting. Additionally, superconductor vortices with continuous chemical potential prevent access to edge-modes between non-adjacent topological phases as was observed for Gaussian vortices. \\
Furthermore, contrary to the prediction that reducing the vortex radius should remove all states present in the superconducting gap, we observed states in the gap that were unaffected by the vortex size as well as states that where removed by small enough vortex sizes.

As topological quantum computing relies on the non abelian exchange of such Majorana modes, the exchange statistics of Majorana modes must persist, even when a small instance of energy-splitting is present. We succeeded in showing, the exchange statistic of two extended vortices on a finite lattice with limited distance and minor energy splitting hold up to a phase in the adiabatic limit. However, a phase common to both the ground and excited state can change a clockwise $T(\sigma_1)$  to counter-clockwise $T(\sigma_1)^{-1}$  exchange and reverse. We demonstrate that, in practice, a dynamical phase from any energy splitting can, unless compensated,  limit access to the adiabatic limit.

We further provided Plots of the Fidelity for the exchange in finite time. Here we saw confirmation that the vortex-to-vortex distance has a larger impact on the fidelity than the vortex radius. We observed the onset of the adiabatic regime at $\omega = 10^{-4} [\frac{\xi}{\hbar 2 \pi}]$  for the exchange of two Gaussian vortices $\mu = 1, \Delta = 1, t = -\frac{2}{3}, R = \frac{1}{35}$.

While the computed grid sizes are comparatively small with $36 \times 36$ sites compared to Abrikosov vortices in type-II superconductors, the vortex size is reasonable. In BCS theory, the vortex core has a radius similar to the superconducting coherence length $\xi$, whereas the vortex-to-vortex distance is of at least the order of the London penetration depth $\lambda$. For YbCuO with lattice constants $a = 3.82, b = 3.89,$ and $c = 11.68 \si{\angstrom}.$ the coherence length in the ab-plane is $\xi_{ab}  \sim 20 \si{\angstrom}.$  This is equivalent to $5.2$ lattice sites when observed directly, the vortices have radii between $20$ and $100 \si{\angstrom}$ \cite{sonier2004investigations} depending on the magnetic field. Thus, the vortex radius of $r = 4.5$ used in our calculations is very similar. However, the London penetration depth in the ab-plane of YBCuO is $\lambda_{ab} \sim150$ nm, equivalent to $375$ lattice sites and $\sim20$ times larger than the $18$ sites used in the simulation. We can expect the exponential energy splitting in dependence of the vortex to vortex/bulk edge to improve the fidelity in a real-world experiment, especially as it was much stronger than the radial energy-splitting.\\
In future calculations, it would be interesting to increase the grid size. A grid of up to $100\times100$ should be feasible with the provided implementation, given enough computing time. Furthermore, one could calculate the fidelity for the standard quantum gates Pauli-X gate $(\tau(\sigma_2)\tau(\sigma_2))$, Hadamard $(\tau(\sigma_1)\tau(\sigma_2)\tau(\sigma_1))$ and CNOT. For Pauli-X and Hadamard, the provided four vortices suffice, while for CNOT, six vortices are required.
\end{document}